\documentclass[10pt,a4paper]{article}

\usepackage{fontspec}
\usepackage[ukrainian]{babel}
\usepackage{geometry}
\usepackage{enumitem}
\usepackage{array}
\usepackage{xcolor}
\usepackage{amsmath}
\usepackage{amssymb}

\setmainfont{Arial}[
  Path=/System/Library/Fonts/Supplemental/,
  UprightFont=Arial.ttf,
  BoldFont=Arial Bold.ttf,
  ItalicFont=Arial Italic.ttf,
  BoldItalicFont=Arial Bold Italic.ttf
]
\geometry{top=0.8cm,bottom=0.8cm,left=1.2cm,right=1.2cm}
\setlength{\parindent}{0pt}
\setlength{\parskip}{2pt}
\setlist[enumerate]{leftmargin=*,itemsep=2pt,topsep=2pt}
\pagestyle{empty}

\newcommand{\answerline}[1]{\par\vspace{2pt}\rule{#1}{0.4pt}}
\newcommand{\longanswer}{\par\vspace{5pt}\rule{\linewidth}{0.4pt}}

\begin{document}

\begin{center}
  {\Large\bfseries Діагностична самостійна робота}\par
  {\large STEM, 6 клас}\par
  \small Робота не оцінюється. Виконуй завдання самостійно й показуй обчислення.
\end{center}

\begin{center}
  {\large\bfseries Варіант 1}
\end{center}

\textbf{Завдання 1. Обчислення}

Виконай дії.
\begin{enumerate}[label=\alph*)]
  \item $48+27$
  \item $90-36$
  \item $6\cdot 8$
  \item $3{,}6+1{,}45$
\end{enumerate}

Виконай письмово (у стовпчик): д) $324\cdot27$;\quad е) $936:24$.

\textbf{Завдання 2. Робота з результатами вимірювання}

Учні вимірювали висоту паростка протягом чотирьох тижнів.

\begin{center}
\begin{tabular}{|c|c|c|c|c|}
\hline
Тиждень & 1 & 2 & 3 & 4 \\
\hline
Висота, см & 8 & 11 & 15 & 18 \\
\hline
\end{tabular}
\end{center}

На скільки сантиметрів виріс паросток від першого до четвертого тижня? Покажи обчислення.
\answerline{\linewidth}

Якою була середня висота паростка? Обчисли середнє арифметичне чотирьох вимірювань.
\answerline{\linewidth}

\textbf{Завдання 3. Знаки та символи}

Познач природний знак літерою П, а створений людиною знак або символ --- літерою Ш:

темні хмари перед дощем --- \rule{0.8cm}{0.4pt};\quad дорожній знак --- \rule{0.8cm}{0.4pt};\quad
слід тварини на снігу --- \rule{0.8cm}{0.4pt};\quad знак «$+$» --- \rule{0.8cm}{0.4pt}.

Обери один приклад і поясни, яку інформацію він передає.
\answerline{\linewidth}

\textbf{Завдання 4. Що входить до STEM}

Установи відповідність. Біля кожної літери запиши номер правильного пояснення.

\begin{tabular}{@{}p{0.43\linewidth}p{0.53\linewidth}@{}}
S --- \rule{1cm}{0.4pt} & 1. Проєктування і створення конструкцій\\[3pt]
T --- \rule{1cm}{0.4pt} & 2. Обчислення, вимірювання та закономірності\\[3pt]
E --- \rule{1cm}{0.4pt} & 3. Дослідження природи та її явищ\\[3pt]
M --- \rule{1cm}{0.4pt} & 4. Інструменти, пристрої та цифрові рішення\\[3pt]
& 5. Механічне запам'ятовування тексту без дослідження
\end{tabular}

\textbf{Завдання 5. Плануємо STEM-проєкт}

У класі швидко висихає ґрунт у вазонах. Учні хочуть створити простий індикатор вологості.

\begin{enumerate}[label=\alph*)]
  \item Сформулюй проблему одним реченням.
  \answerline{\linewidth}
  \item Сформулюй мету проєкту.
  \answerline{\linewidth}
  \item Розташуй етапи у правильній послідовності, записавши номери:

  \begin{center}
  випробування моделі --- пошук інформації --- презентація результату --- створення моделі
  \end{center}
  Відповідь: \rule{6cm}{0.4pt}
  \item Назви одну величину, яку потрібно вимірювати або спостерігати під час випробування.
  \answerline{\linewidth}
\end{enumerate}

\newpage

\begin{center}
  {\Large\bfseries Діагностична самостійна робота}\par
  {\large STEM, 6 клас}\par
  \small Робота не оцінюється. Виконуй завдання самостійно й показуй обчислення.
\end{center}

\begin{center}
  {\large\bfseries Варіант 2}
\end{center}

\textbf{Завдання 1. Обчислення}

Виконай дії.
\begin{enumerate}[label=\alph*)]
  \item $37+46$
  \item $80-27$
  \item $7\cdot 9$
  \item $2{,}8+1{,}35$
\end{enumerate}

Виконай письмово (у стовпчик): д) $418\cdot35$;\quad е) $864:27$.

\textbf{Завдання 2. Робота з результатами вимірювання}

Учні вимірювали температуру води під час досліду.

\begin{center}
\begin{tabular}{|c|c|c|c|c|}
\hline
Вимірювання & 1 & 2 & 3 & 4 \\
\hline
Температура, $^\circ$C & 18 & 22 & 26 & 30 \\
\hline
\end{tabular}
\end{center}

На скільки градусів змінилася температура від першого до четвертого вимірювання? Покажи обчислення.
\answerline{\linewidth}

Якою була середня температура? Обчисли середнє арифметичне чотирьох вимірювань.
\answerline{\linewidth}

\textbf{Завдання 3. Знаки та символи}

Познач природний знак літерою П, а створений людиною знак або символ --- літерою Ш:

іній на траві --- \rule{0.8cm}{0.4pt};\quad шкільний дзвінок --- \rule{0.8cm}{0.4pt};\quad
блискавка --- \rule{0.8cm}{0.4pt};\quad літера «А» --- \rule{0.8cm}{0.4pt}.

Обери один приклад і поясни, яку інформацію він передає.
\answerline{\linewidth}

\textbf{Завдання 4. Що входить до STEM}

Установи відповідність. Біля кожної літери запиши номер правильного прикладу діяльності.

\begin{tabular}{@{}p{0.43\linewidth}p{0.53\linewidth}@{}}
S --- \rule{1cm}{0.4pt} & 1. Створити креслення мосту та його модель\\[3pt]
T --- \rule{1cm}{0.4pt} & 2. Обчислити витрати матеріалів\\[3pt]
E --- \rule{1cm}{0.4pt} & 3. Дослідити умови проростання насіння\\[3pt]
M --- \rule{1cm}{0.4pt} & 4. Запрограмувати датчик температури\\[3pt]
& 5. Переписати умову завдання, не виконуючи дослідження
\end{tabular}

\textbf{Завдання 5. Плануємо STEM-проєкт}

У шкільному коридорі світло залишається ввімкненим навіть тоді, коли там нікого немає. Учні хочуть створити модель автоматичного освітлення.

\begin{enumerate}[label=\alph*)]
  \item Сформулюй проблему одним реченням.
  \answerline{\linewidth}
  \item Сформулюй мету проєкту.
  \answerline{\linewidth}
  \item Розташуй етапи у правильній послідовності, записавши номери:

  \begin{center}
  презентація результату --- створення моделі --- пошук інформації --- випробування моделі
  \end{center}
  Відповідь: \rule{6cm}{0.4pt}
  \item Назви одну величину, яку потрібно вимірювати або спостерігати під час випробування.
  \answerline{\linewidth}
\end{enumerate}

\end{document}
